\documentclass[12pt]{article}

% Packages
\usepackage{graphicx}      % For logo
\usepackage{tcolorbox}     % For note box
\usepackage{setspace}      % Optional spacing
\usepackage{geometry}      % Page layout
\usepackage{draftwatermark}

\geometry{margin=1in}

\begin{document}

% Watermark:
\SetWatermarkText{DRAFT} % Text of the watermark
\SetWatermarkScale{4} % Size (scale factor)
\SetWatermarkColor[gray]{0.85} % Transparency (0 = invisible, 1 = fully opaque)

% -------------------------------
% TITLE PAGE
% -------------------------------
\begin{titlepage}
    \centering

    % Top row: Logo (left) and Doc Info (right)
    \noindent
    \begin{minipage}[t]{0.49\textwidth}
        \vspace{0pt}
        \includegraphics[width=0.6\textwidth]{logo_text.png}
    \end{minipage}
    \hfill
    \begin{minipage}[t]{0.49\textwidth}
        \vspace{0pt}
        \raggedleft
        \begin{tabular}{rl}
            \textbf{Document Number:} & RANF-ENG-00004 \\
            \textbf{Revision:} & Rev. DRAFT \\
        \end{tabular}
    \end{minipage}

    % Vertical spacing before title
    \vspace{2.5cm}

%    % Document Number and Revision
%    \begin{tabular}{rl}
%        \textbf{Document Number:} & DOC-00123 \\
%        \textbf{Revision:} & Rev. A \\
%    \end{tabular}
%
%    % Organization Logo
%    \vspace*{1cm}
%    \includegraphics[width=0.35\textwidth]{ranf_imprint_logo.png}\\[1.5cm]

% TODO: add QR code to website

    % Report Title
    {\Huge \bfseries Final / Critical Design Review \\[0.5cm]}

    % Author and Date
    {\Large Timothy P Woodard \\[0.2cm]}
    {\large \today \\[1.5cm]}

    % Note Box
    \begin{tcolorbox}[colback=gray!10, colframe=black, width=0.9\textwidth]
        \textbf{Note:} Copyright 2026 Roads Are No Fun

This document is provided as part of the Roads Are No Fun project under the Creative Commons CC BY-SA license.
    \end{tcolorbox}

    \vfill
\end{titlepage}

% -------------------------------
% MAIN DOCUMENT
% -------------------------------
%\tableofcontents
\newpage

\section{Introduction}
Design reviews serve as gating events that establish the progress attained during the design process for a product and determine whether the next steps of design/production can begin, or if additional work is required.
The final or critical design review is the review which evaluates if the design is complete and ready for production.
Obviously at this stage of design, all requirements and constraints must be expected to be met, or have a clear path identified to be met.

\section{Expectations}

\subsection{Front Matter}

point of contact / author / organization

date, revision 

subject / identifiers / part numbers

contents


\subsection{Design Overview}

context

purpose 

function

overall design concept / architecture

consists of


\subsection{Design Criteria}

guidance documents / standards

requirements:

performance criteria

functional criteria

requirements verification matrix

key performance parameters

higher level driving requirements

derived requirements

concept specific requirements


\subsection{Design Details}

detailed architecture

what and how

standard / COTS components

materials


trade studies / alternatives / evaluations


\subsection{Lifecycle Summary}



manufacture

processes

assembly

operation

maintenance

disposal

tooling and test equipment

facilities


\subsection{Design Verification}

verification plan: inspection, analysis, test, demonstrate, similarity, lot testing

qualification and acceptance

Analysis
-methodology
-margins summary
-loads and environments
-analysis results

durability: fracture / fatigue / wear

Testing: test plan


\subsection{Safety \& Failure}

FMECA

Hazards

Hazard controls

risks and watch items


\subsection{Forward Work}

Identify and remaining open items

schedule and cost estimates


\subsection{Success / Exit Criteria}

design is complete

design meets all requirements and constraints

feasible to build

ready for implementation / production





%\section{Glossary}
%\paragraph{Not Invented Here (NIH) Syndrome} when engineers at a company do not use a system external to their organization due to the perceived inferiority to internally developed systems
%\paragraph{C-Sketch} collaborative sketching, also called 6-3-5, is a brainstorming method

\section*{References}

\begin{enumerate}
\item MIL-A-83577, 1988.
    %\item George E. Dieter, Linda C. Schmidt; Engineering Design Fourth Edition; McGraw Hill; 2009.
\end{enumerate}

\end{document}